% AUTOMATISCH ERZEUGT aus dem SysML-Modell mit tools/generate_docs.py.
% Nicht von Hand aendern: Inhalte im Modell pflegen, Layout im Kapitel.
SN-OPS-001 & STK-003 & Der Betreiber moechte den Betriebszustand des Roboters erkennen koennen. \\
SN-OPS-002 & STK-003, STK-005 & Stoerungen sollen nachvollziehbar sein. \\
SN-OPS-003 & STK-003 & Der Roboter soll im normalen Betrieb moeglichst wenig manuelle Eingriffe benoetigen. \\
SN-MNT-001 & STK-005 & Wartungspersonal soll Fehler diagnostizieren koennen. \\
SN-MNT-002 & STK-005 & Verschleissteile und zentrale Komponenten sollen zugaenglich sein. \\
SN-MNT-003 & STK-005, STK-007 & Hardware und Software sollen austauschbar beziehungsweise aktualisierbar sein. \\
SN-DEV-001 & STK-009 & Das System soll schrittweise entwickelt und erweitert werden koennen. \\
SN-DEV-002 & STK-009 & Aenderungen sollen moeglichst lokal bleiben und keine unnoetigen Aenderungen anderer Subsysteme verursachen. \\
SN-DEV-003 & STK-009, STK-010 & Anforderungen, Architektur, Implementierung und Tests sollen nachvollziehbar miteinander verbunden sein. \\
